% Background: Sources F1--F7 (equilibrium + game theory convergence);
% cites McKelvey1995, Sokota2023, Mertikopoulos2018, Perolat2021, Loshchilov2019,
% Bai2019, Anderson1965, Broyden1965, Geng2021, Winston2020, Miyato2018, Dehghani2019,
% Lan2020, Graves2016, RecursiveTransformers2024, Lee2015, Rafailov2023, Wu2024, Wang2025MPO,
% Duetting2024, Vickrey1961, Holtzman2020, Bengio2015, Warstadt2020, Warstadt2023, Radford2019.
% No invented numbers. Drawn from ADR 0006 (PMA vs MPO distinction).

\section{Background and Related Work}
\label{sec:background}

Contemporary language model optimization is dominated by a single loss function and a single optimizer.
One might caricature this regime as \emph{dictatorial optimization}.
Yet the systems in which these models operate are inherently interactive: preference tuning optimizes
a policy against a frozen reference, decoding aggregates diverse model outputs, and training increasingly
unfolds as a feedback loop between model, data, and reference.
Game-theoretic machinery offers the natural mathematical language for such settings.
A growing line of recent work has begun to import equilibrium theory and mechanism design into language modeling
\citep{Sokota2023, Bai2019, Duetting2024, Wu2024}.
This section reviews the foundational concepts and prior work that motivate the equilibrium-theoretic treatment
of three stages of language modeling: training dynamics via magnetic mirror descent, model depth via implicit
fixed-point computation, and decoding via truthful mechanism design.

\subsection{Quantal Response Equilibrium and Softmax as Bounded Rationality}

Quantal response equilibrium (QRE), introduced by McKelvey and Palfrey in a foundational 1995 paper on
normal-form games, models boundedly rational agents by replacing best-response play with stochastic best
response \citep{McKelvey1995}.
Under a QRE, an agent plays action $a$ with probability proportional to $\exp(\lambda\, u(a))$,
where $u(a)$ is the payoff and $\lambda$ is a rationality (inverse-temperature) parameter that captures
the agent's degree of strategic optimization.
As $\lambda \to \infty$, the QRE converges to Nash equilibrium; at intermediate $\lambda$ it interpolates
between uniform random play (low $\lambda$) and best response (high $\lambda$), and coincides with
entropy-regularized equilibrium.
The normalizing constant, $\sum_a \exp(\lambda\, u(a))$, is the softmax partition function.
QRE has been widely studied in experimental economics and behavioral game theory to characterize how real
players deviate from perfect rationality.

The connection to language modeling is direct and structural.
The softmax output layer of any modern language model is literally a quantal response to its logits.
At each token position, the model computes a vector of scores (logits $\ell \in \mathbb{R}^V$) and outputs
a probability distribution via softmax: $p_i = \exp(\ell_i) / Z$, where $Z = \sum_j \exp(\ell_j)$ is the
partition function.
The temperature parameter $\tau = 1/\lambda$ controls the peakedness of this distribution.
Low temperatures (high $\lambda$) concentrate probability sharply on high-scoring tokens; high temperatures
spread probability more evenly across continuations.
This observation, though elementary, motivates a systematic study of language model decoding through the
lens of QRE: the model's token distribution traces out QRE paths as temperature varies, and one may
interrogate what equilibrium properties the learned logits satisfy.
Such paths arise naturally in games where agents adapt their strategies in response to fixed payoff structures.
Understanding whether language models' learned representations align with QRE predictions is an open empirical
question and is explored in Section~\ref{sec:convergence} of the present work.

\subsection{Mirror Descent, Optimization Geometry, and Regularized Learning}

Mirror descent, a classical technique in convex optimization, generalizes gradient descent by replacing the
Euclidean projection step with a projection in a non-Euclidean geometry defined by a \emph{mirror map}
(or Bregman divergence).
When applied to probability distributions on a simplex, using the negative entropy as the mirror map naturally
yields softmax updates, making it well-suited to probability-valued optimization.
The convergence rate and convergence mode (oscillatory vs.~monotone decay) depend on both the choice of
mirror map and on the learning rate.
Mirror descent is known to achieve faster convergence than vanilla gradient descent on certain problem
classes and is particularly effective for saddle-point problems in two-player games where it avoids cycling.

\subsection{Magnetic Mirror Descent and Convergence in Two-Player Games}

Magnetic Mirror Descent (MMD), introduced by Sokota et al.\ at ICLR 2023, augments standard mirror descent
with a proximal pull toward a reference (magnet) policy $y_{\mathrm{ref}}$ \citep{Sokota2023}.
Using the negative-entropy mirror map and regularization strength $\tau$, the update admits a closed form:
\begin{equation}
y_{t+1} \;=\; \frac{y_t + \eta g_t + \eta\tau\, y_{\mathrm{ref}}}{1 + \eta\tau},
\label{eq:mmd}
\end{equation}
where $g_t$ is the payoff gradient and $\eta$ the step size.
Sokota et al.\ prove linear (geometric) last-iterate convergence to the $\tau$-regularized equilibrium in
two-player zero-sum games.
This property is significant because it contrasts sharply with the behavior of simultaneous gradient
descent-ascent (GDA), which cycles away from equilibrium under identical problem structure and step sizes.
Mertikopoulos et al.\ established this cycling behavior in foundational work \citep{Mertikopoulos2018}.
The magnetic term thus breaks the cycling pathology that plagues standard GDA, offering convergence guarantees
in settings where unregularized gradient play fails to approximate equilibrium.

Annealing the regularization away permits recovery of Nash equilibrium itself, not merely the
$\tau$-regularized variant.
Periodically re-centering the magnet on the current iterate anneals $\tau$ toward zero.
This technique is rooted in regularized Nash dynamics \citep{Perolat2021}, which shows that re-centering on
the running average recovers Nash convergence in both symmetric and asymmetric games.
The parameter-space variant, MagneticAdamW, applies the magnetic pull after a decoupled-weight-decay AdamW
step \citep{Loshchilov2019}, with the reference policy either an exponential moving average of recent iterates,
a periodic snapshot, or frozen base weights (as employed in preference optimization).
The closed-form update of Eq.~\ref{eq:mmd} carries directly into weight space, replacing the policy $y$ with
weight parameters and $g_t$ with the normalized gradient direction from adaptive moment estimation.

\subsection{Deep Equilibrium Models: Implicit Depth and Fixed-Point Computation}

Deep equilibrium models (DEQ), proposed by Bai et al.\ at NeurIPS 2019, replace a traditional $L$-layer
feed-forward stack with a single parameterized map $f_\theta$ solved to a fixed point \citep{Bai2019}.
The forward pass solves $z^* = f_\theta(z^*, x)$ iteratively until convergence; the representation $z^*$ is
extracted and used for downstream prediction.
The input $x$ is fed in at every iteration, so the model's output depends on both the equilibrium
representation and the input conditioning.
Gradients flow backward through the fixed-point computation using the implicit function theorem, which avoids
unrolling the solver and storing intermediate activations during the forward pass.

The activation memory advantage is substantial and measurable.
For an $L$-layer explicit stack with activation dimension $d$, storing activations for backpropagation costs
$O(Ld^2)$ memory; a DEQ model with the same dimension and iteration budget stores only $O(d^2)$ activation
memory, independent of effective depth.
For a 12-layer explicit stack with $d=768$, the savings are dramatic: $O(12 \cdot 768^2)$ vs.~$O(768^2)$.
This memory advantage becomes critical at large effective depths or when computing backward passes on
limited-memory devices, making DEQ particularly attractive for memory-constrained settings.

Numerically, the fixed point is computed via Picard iteration (repeatedly applying the map), Anderson
acceleration \citep{Anderson1965}, or Broyden's quasi-Newton method \citep{Broyden1965}.
Picard iteration is simple but can be slow, especially for poorly-conditioned fixed-point equations.
Anderson acceleration uses a low-rank history (typically 3--5 prior iterates) to construct a linear
interpolant that accelerates convergence and approaches quasi-Newton efficiency on many problem classes.
Broyden's method, a quasi-Newton variant, maintains a low-rank approximation to the Jacobian and can be faster
still on well-scaled problems.
Backward pass computation can use the full implicit differentiation formula (solving a linear system at
gradient time) or cheaper Jacobian-free approximations \citep{Geng2021}.

A critical requirement is that the fixed-point map must be contractive, or at minimum converge to a fixed
point.
Spectral normalization \citep{Miyato2018}, borrowed from generative adversarial networks, constrains the
Lipschitz constant of the map; parametrizing layers as $W / \|W\|_2$ (or using iterative power iteration)
ensures Lipschitz constant $\leq 1$ and guarantees contraction.
Alternatively, monotone operator networks \citep{Winston2020} restrict the map to a class that admits a unique
fixed point by architectural design.
Failure to enforce contraction results in diverging solvers and non-existent fixed points, a practical pitfall
encountered when designing equilibrium language models and is discussed in Section~\ref{sec:diagnostic}.

\subsection{Weight Tying, Recursive Uptraining, and Deep Supervision}

Weight tying---reusing identical parameters across the $L$ iterations or depth of a deep model---is
independently motivated by parameter efficiency and is orthogonal to the fixed-point formulation itself.
Explicit architectural examples include Universal Transformers \citep{Dehghani2019}, which tie weights across
time steps and add learned per-instance halting.
ALBERT \citep{Lan2020} shares embeddings across layers and intermediate feed-forward layers to reduce parameter
count.
Adaptive Computation Time \citep{Graves2016} learns per-instance step budgets dynamically.
Deep supervision \citep{Lee2015} provides auxiliary losses at intermediate layers during training, giving
gradient signals at each iteration rather than only at the final output; this has been shown empirically to
help with optimization and generalization in deep networks.
Recursive uptraining \citep{RecursiveTransformers2024} fine-tunes pretrained models via repeated application
of learnable layer-wise updates, permitting adaptation of models trained at a fixed depth to operate at
multiple effective depths with shared parameters.
These techniques are complementary to DEQ and can be combined to obtain benefits of both parameter sharing and
implicit depth.

\subsection{Preference Optimization: Mechanisms and the Critical Distinction Between Policy and Parameter Space}

Direct Preference Optimization (DPO) \citep{Rafailov2023}, introduced by Rafailov et al.\ at NeurIPS 2023,
optimizes the language model directly against a preference signal without training a separate reward model.
DPO frames alignment as a maximum-likelihood problem with a closed-form solution in terms of the reference
model's log-probabilities, making it simpler and often more stable than methods requiring auxiliary reward
models.
The loss is typically $-\log \sigma(\lambda(s_w - s_l))$, where $s_w$ and $s_l$ are the log-probabilities
assigned by the model to preferred and dispreferred completions, and $\lambda$ is a temperature parameter
controlling the strength of the preference signal.

Policy-space Magnetic Preference Optimization (MPO), published by Wang et al.\ in the ICLR 2025 cycle,
augments DPO with magnetic mirror descent in the policy (logit) space \citep{Wang2025MPO}.
MPO applies the proximal magnetic term to the model's logits while optimizing the standard DPO objective
under self-play, where the model's own generations form the preference data for training.
The Nash equilibrium of the resulting preference game becomes the target of the magnetic pull in policy space.

This is structurally different from parameter-space magnetic anchoring (PMA), the method studied in the present
work.
PMA applies a proximal pull in the \emph{parameter} space toward frozen base weights.
The magnetic term may be folded into the optimizer---by scaling the Adam update direction and mixing in a pull
toward the frozen reference---or applied outside the optimizer as a decoupled exponential-moving-average drift.
Critically, PMA operates on the standard DPO loss function without self-play or game-theoretic restructuring of
the preference data.
Though both policy-space MPO and parameter-space PMA deploy magnetic terms and aim to improve preference
alignment, they operate in fundamentally different representation spaces (logits vs.~parameters), with
different target equilibria (Nash in preference space vs.~frozen reference), and under different game-theoretic
framings (self-play vs.~anchor-relative optimization).
Consequently, results concerning the effectiveness or ineffectiveness of PMA carry no direct bearing on
policy-space MPO and should not be conflated.
Self-Play Preference Optimization (SPPO) \citep{Wu2024} pursues alignment through iterative self-play without
an explicit magnetic term, offering yet another orthogonal mechanism.

\subsection{Mechanism Design and Truthful Aggregation of Multiple Models}

When multiple models jointly produce a single output, each model ``bids'' its confidence in a continuation,
and an aggregation mechanism decides which model to select---a canonical mechanism-design problem
\citep{Duetting2024}.
A mechanism is truthful (incentive-compatible or strategy-proof) if each agent's best response is to
truthfully report its private information; this property prevents gaming and ensures that aggregation reflects
genuine model confidence rather than strategic misreporting to gain influence.

Vickrey's second-price auction \citep{Vickrey1961}, developed in the context of sealed-bid auctions, is the
canonical truthful mechanism for single-item auctions.
In the context of token-level decoding, the ``item'' is the token position, each model's ``valuation'' is its
predicted confidence (log-probability or entropy), and the second-price rule (the winner pays the
second-highest bid) ensures that truthful bidding is a dominant strategy.
A model cannot profit by bidding higher (it wins but pays more) or lower (it loses its opportunity to
contribute).
Alternative aggregation schemes, such as weighted averaging of logits or softmax outputs, are not truthful;
models face incentives to misreport their confidence to increase their influence on the aggregated output.

Beyond truthfulness, practitioners must contend with text degeneration, where greedy or even uniform
aggregation of model outputs can produce repetitive or incoherent text \citep{Holtzman2020}.
Exposure bias in autoregressive generation \citep{Bengio2015} refers to the distribution mismatch between
training (where the model conditions on reference data) and inference (where it conditions on its own outputs);
scheduled sampling is one approach to mitigate this.
Auction mechanisms do not directly solve exposure bias but provide a principled framework for multimodel
decoding grounded in mechanism theory and economic incentives.

\subsection{Language Model Benchmarks and Sample-Efficient Evaluation}

The Benchmark of Linguistic Minimal Pairs for English (BLiMP) evaluates grammatical and semantic acceptability
through contrastive pairs across a range of phenomena---agreement, binding, governing category, and others
\citep{Warstadt2020}.
BLiMP has become the standard for assessing small language models' linguistic competence without requiring
large supervised instruction-tuned datasets; it provides a reliable signal at the 10M--1B parameter scale.
The BabyLM Challenge \citep{Warstadt2023} established a shared evaluation protocol for models pretrained on a
fixed small corpus (10 million words, divided into strict-small, strict-medium, and strict-large
configurations), enabling direct and controlled sample-efficient comparisons across methods.
The benchmark family is critical for validating methods at pre-scaling scales, where GPU compute is limited
and researchers can iterate quickly on architectural and training innovations.
Separately, GPT-2's evaluation on standard perplexity-based benchmarks \citep{Radford2019} provides reference
points for autoregressive modeling quality on diverse text corpora including Common Crawl and WebText, though
at larger scales where pretraining dominates.

\subsection{Equilibrium, Depth, and Certification: A Comparative Landscape}

\begin{table}[H]
\centering
\caption{Equilibrium-theoretic and depth-implicit architectures: parameter sharing, fixed-point semantics,
convergence guarantees, and certification status. No single architecture achieves certified convergence while
preserving linguistic quality at the BabyLM scale in the present work.}
\label{tab:equil_comparison}
{\small
\begin{tabular}{llcccc}
\toprule
\textbf{Framework} & \textbf{Param.\ Sharing} & \textbf{Depth Semantics} & \textbf{Fixed Pt.} & \textbf{Certified} \\
\midrule
Explicit stacking & none & sequential & N/A & N/A \\
Universal Transformer & tied weights & learned halting & conditional & per-instance \\
ALBERT & tied embeddings & parameter reduction & N/A & N/A \\
Deep Equilibrium (DEQ) & single block & fixed-point equation & yes & via spectral norm \\
Monotone Operator Nets & constrained weights & fixed-point enforced & yes & by construction \\
EqLM (anytime unroll) & tied block & supervised iteration & trained & no \\
EqLM (certified) & tied block & fixed-point w/ penalty & yes, locally & trajectory-local \\
\bottomrule
\end{tabular}
}
\end{table}

The table organizes the landscape of depth-implicit and equilibrium approaches along four axes.
Parameter sharing distinguishes methods that reuse parameters (reducing total model size) from those that stack
unique layers independently.
Depth semantics characterizes how the framework interprets ``depth'': sequential composition, learned
per-instance adaptation, fixed-point computation, or structural constraint enforcement.
Fixed point indicates whether the architecture provably admits a well-defined equilibrium state or fixed point.
Certification records whether convergence to that equilibrium is guaranteed by architectural design or training
regularization.
The row for EqLM's certified variant (Section~\ref{sec:h6}) reflects a key empirical finding: a
trajectory-local finite-difference penalty enforces convergence, yielding verifiable fixed-point correctness at
the 121M parameter scale, but at the cost of linguistic quality (BLiMP score degradation).
This leaves quality-preserving certification as a sharply posed open problem at the present scale and
architecture, motivating the scale-up program described in Section~\ref{sec:method}.
